\begin{proposition}\label{rl-0018}\label{prop:compatible-triples-for-main-proof}
  In the above notation,
  \begin{itemize}[(i)]
    \item \(\check \phi\) is a \(T\)-equivariant relative perfect obstruction theory for \(\pi^T:M^T\to \mathcal Z\)
    \item (First triple) ~ \((j^*\widetilde{\mathbb E)}, \check {\mathbb E}, N^\vee_0[1]\) is a compatible triple.
    \item (Second Triple) ~ \((\mathbb E^f, \check{\mathbb E}, N_1^\vee[1])\) is a compatible triple.
  \end{itemize}
\end{proposition}
\begin{proof}
$ $

  \noindent \emph{(i)} \quad \(\check{\mathbb E}\) is the direct sum of two complexes each of perfect amplitude \([-1,0]\), so it is of perfect amplitude \([-1,0]\) as well. Since \(h^0(N^\vee_1[1]) = 0\) we get \(h^0(\check \phi) = h^0(\phi^f)\) and \(h^{-1}(\check \phi)\) restricted to the summand \(h^{-1}(\mathbb E^f)\) is \(h^{-1}(\phi^f)\), hence \(h^0(\check \phi)\) and \(h^{-1}(\check \phi)\) are an isomorphism and a surjection respectively by \ref{lem:obstruction-theory-for-the-fixed-locus}.

  \bigskip

  \noindent \emph{(ii)} \quad The diagram for the first compatible triple is
  \begin{center}
    \begin{tikzcd}[row sep = large, column sep = large]
      j^*\widetilde{\mathbb E} = \mathbb E^f\oplus \mathbb E^m
        \arrow[r, "\beta"]
        \arrow[d, "j^*\phi"']
      & \check{\mathbb{E}} = \mathbb{E}^f \oplus N_1^\vee[1]
        \arrow[r, "\check\rho"]
        \arrow[d, "\check\phi"]
      & N_0^\vee[1]
        \arrow[r, "+1"]
        \arrow[d, "\mathfrak f"]
      & {} \\
      j^* L_{\widetilde{M}/\mathcal{Y}^T}
        \arrow[r]
      & L_{M^T/\mathcal{Y}^T}
        \arrow[r]
      & L_{M^T/\widetilde{M}}
        \arrow[r, "+1"]
      & {}
    \end{tikzcd}
  \end{center}
  Here
  \begin{align*}
    \beta = \id_{\mathbb E^f}\oplus q, \quad q:[N_1^\vee \xrightarrow{d^\vee} N^\vee_0] \twoheadrightarrow [N^\vee_{1}\to 0] = N^\vee_1[1]
  \end{align*}
  and \(\check \rho = (0,d^\vee [1])\).

  The top row is a compatible triple because it is the direct sum of the compatible triples \(\mathbb E^F\xrightarrow{\id_{\mathbb E^f}} \mathbb E^f \xrightarrow{0}0\) and \(\mathbb E^m \xrightarrow{q} N_1^\vee[1] \xrightarrow{d^\vee[1]} N_1^\vee[1]\).

  The map \(\mathfrak f\) exists by axiom (TR3) completing the two rows to a morphism of triangles. Fix any such map. We are then in exactly the situation of \cite[Construction 3.13]{manolache_VirtualPullbacks2012} and the diagram chase there shows \(\mathfrak f\) is a perfect obstruction theory.

  \bigskip
  
  \noindent \emph{(iii)} \quad The diagram for the second compatible triple is
  \begin{center}
    \begin{tikzcd}[row sep = large, column sep = large]
      \mathbb{E}^f
        \arrow[r, "{(\mathrm{id},\,0)}"]
        \arrow[d, "\phi^f"']
      & \check{\mathbb{E}} = \mathbb{E}^f \oplus N_1^\vee[1]
        \arrow[r, "{(0,\,\mathrm{id})}"]
        \arrow[d, "\check\phi"]
      & N_1^\vee[1]
        \arrow[r, "0"]
        \arrow[d, "0"]
      & {} \\
      L_{M^T/\mathcal{Y}^T}
        \arrow[r, "\mathrm{id}"]
      & L_{M^T/\mathcal{Y}^T}
        \arrow[r]
      & L_{\mathrm{id}_{M^T}} = 0
        \arrow[r, "+1"]
      & {}
    \end{tikzcd}
  \end{center}
  Checking that each vertical map is a perfect obstruction theory and that the whole thing commutes is routine.
\end{proof}
